\section{Preliminaries}\label{sec:prelims}

\subsection{Definitions} 

% \begin{definition}[Key Encapsulation Mechanisms]
% A Key Encapsulation Mechanism (KEM) is a tuple of algorithms $(\mathsf{KeyGen}, \mathsf{Encap}, \mathsf{Decap})$ defined as follows:
% \begin{itemize}
%     \item $(\mathsf{pk}, \mathsf{sk}) \leftarrow \mathsf{KeyGen}()$: generates a public/private key pair.
%     \item $(\mathsf{ct}, \mathsf{ss}) \leftarrow \mathsf{Encap}(\mathsf{pk})$: given a public key, outputs a ciphertext $\mathsf{ct}$ and a shared secret $\mathsf{ss}$.
%     \item $\mathsf{ss} \leftarrow \mathsf{Decap}(\mathsf{sk}, \mathsf{ct})$: given a secret key and ciphertext, recovers the shared secret.
% \end{itemize}
% A KEM is correct if decapsulation recovers the same shared secret produced by encapsulation almost always, and secure if it satisfies IND-CCA (indistinguishability under adaptive chosen-ciphertext attack)~\cite{EPRINT:Shoup01}.
% \end{definition}

% During the current transition period, hybrid KEMs combine a classical scheme with a post-quantum scheme, where security holds if either component remains unbroken.
% In this work, we instantiate hybrid key exchange using X-Wing~\cite{EPRINT:BCDKSV24}, which combines X25519~\cite{PKC:Bernstein06} and ML-KEM-768~\cite{NIST:24MLDSA} into a single KEM.

\begin{definition}[Digital Signatures]
A digital signature scheme is a tuple of algorithms $(\mathsf{KeyGen}, \mathsf{Sign}, \mathsf{Verify})$ defined as follows:
\begin{itemize}
    \item $(\mathsf{pk}, \mathsf{sk}) \leftarrow \mathsf{KeyGen}()$: generates a public/private key pair.
    \item $\sigma \leftarrow \mathsf{Sign}(\mathsf{sk}, m)$: produces a signature $\sigma$ on message $m$.
    \item $\{0, 1\} \leftarrow \mathsf{Verify}(\mathsf{pk}, \sigma, m)$: returns 1 if the signature is valid, 0 otherwise.
\end{itemize}
A signature scheme is correct if $\mathsf{Verify}(\mathsf{pk}, \mathsf{Sign}(\mathsf{sk}, m), m) = 1$ for all messages $m$ and corresponding keypairs $(\mathsf{pk}, \mathsf{sk})$, and secure if it satisfies EUF-CMA (existential unforgeability under chosen-message attack)~\cite{MenVanVan96}.
\end{definition}

The elliptic curve digital signature algorithm (ECDSA)~\cite{NIST:FIPS1865} is a digital signature scheme based on elliptic curves.
In Ethereum~\cite{ETH:W14}, ECDSA is instantiated over the secp256k1 curve, and a signature on a message $m$ consists of three components: $(\texttt{r}, \texttt{s}, \texttt{v})$, where $\texttt{r}$ and $s$ are 32-byte integers and $\texttt{v}$ is a recovery parameter.
The recovery parameter $\texttt{v}$ encodes information about the parity of the public key $y$-coordinate. 
Therefore, given $(\texttt{r}, \texttt{s}, \texttt{v})$ and $m$, the public key can be uniquely determined without explicitly transmitting it.
This property is exploited in Ethereum during transaction validation~\cite[Appendix F]{ETH:W14}.


\subsection{EVM-based Blockchains}

Blockchains are decentralized systems for maintaining a shared, append-only ledger among a set of untrusted parties. 
Instead of relying on a central authority, a blockchain uses cryptographic techniques, such as hash functions and digital signatures, and distributed consensus to ensure that all parties agree on the order and contents of recorded transactions. 
% Data are organized into blocks that are cryptographically linked, making the ledger tamper-evident once entries are confirmed. 
% This design enables secure, transparent coordination across a network, and has made blockchains a foundational primitive for cryptocurrencies, distributed applications, and other trust-minimized systems.

Ethereum virtual machine (EVM)-based blockchains~\cite{ETH:BU13} organize their protocol stack into distinct layers, each responsible for a specific functionality.
Next, we describe the four primary layers of blockchain systems.

\paragraph{Transactions.}
Transactions are the units of state change in a blockchain.
In Ethereum, transactions use recursive length prefix (RLP) encoding~\cite[Appendix B]{ETH:W14} as the canonical serialization format, providing deterministic encoding.
Since EIP-2718~\cite{eip2718}, Ethereum transactions follow a typed envelope format:
\begin{equation*}
\texttt{Tx} = \texttt{type} \parallel \mathsf{RLP}(\texttt{payload}),
\end{equation*}
where \texttt{type} is a single-byte identifier distinguishing transaction formats (e.g., \texttt{0x02} for EIP-1559 fee market transactions).
This typed structure enables protocol evolution by allowing new transaction formats and we redesign it to accommodate different signature schemes (cf. Section~\ref{s:catx}).

\paragraph{Execution Layer.}
The execution layer (EL) consists of nodes, called clients, which verify and execute transactions and update the network's global state.
This involves executing smart contract bytecode in a virtual machine.
% The execution layer maintains the world state---a mapping from addresses to account states (balance, nonce, code, storage)---and applies state transitions according to transaction semantics.
A specific category of operations in EVM bytecode are precompiled contracts, which offer efficient implementations of computationally intensive cryptographic functions that would otherwise be prohibitively expensive to execute.

\paragraph{Consensus Layer.}
The consensus layer (CL) secures the integrity of block finalization, ensuring all honest nodes agree on the canonical chain of blocks.
Byzantine fault tolerant (BFT)~\cite{ACM:LSP82,ARXIV:BKM18} consensus protocols enable agreement even when some validators behave maliciously.

In weighted BFT systems such as Tendermint~\cite{ARXIV:BKM18}, each validator has an associated voting power and a block is finalized when validators representing more than $2/3$ of voting power sign commit messages for that block.
The consensus proceeds in rounds consisting of three phases: \emph{Propose} (proposer broadcasts a candidate block), \emph{Prevote} (validators vote for the proposal), and \emph{Precommit} (validators commit upon seeing sufficient prevotes).
We use this structure to allow for cryptographic agility in the validator's signature scheme used (cf. Section~\ref{ss:cryptographic-agile-consensus}).

\paragraph{P2P Communication Layer.}
The peer-to-peer (P2P) layer enables nodes to discover each other and exchange messages without centralized coordination.
This layer consists of two main components: a discovery protocol that maintains a distributed hash table (DHT) for peer discovery, and a transport protocol that establishes secure communication channels.

In Ethereum's execution layer, the Discv4 protocol~\cite{ethereum_discv4} enables peer discovery via UDP gossip.
Discovery packets are authenticated using ECDSA signatures with public key recovery, binding messages to node identities.
The RLPx transport protocol~\cite{ethereum_rlpx} then establishes encrypted sessions between discovered peers using ECIES~\cite{JCSE:GHS10} for key agreement.
%Section~\ref{s:p2p-node-communication} presents our modifications to these protocols for post-quantum security.
